\section{Alternatives and Research Position}
Visual authentication is established. Seeing-Is-Believing used two-dimensional
barcodes and camera phones for authentication and demonstrative identification
of devices~\cite{sib}. QRLogin research studies the complete login workflow and
finds failures beyond decoding~\cite{qrlogin}; pQRAuth adds a physical pose
factor~\cite{pqrauth}. OpenAST does not reproduce those deployment studies,
provide their physical factor or claim that signed visual codes are new.

An authenticated web API with transaction-bound permission can already satisfy
many OpenAST requirements. A static signed marker can establish issuer identity
and message integrity but needs additional context for freshness and action
scope. A challenge response improves interaction binding while still admitting
live relay. A pose, distance-bounding or trusted-location mechanism would add
separate evidence, not a stronger interpretation of the same QR signature.

RATS distinguishes evidence, appraisal and relying-party decisions~\cite{rats}.
CaMeL and Fides place capability or information-flow policies around agent
execution~\cite{camel,fides}. OpenAST follows the same architectural discipline;
it is not an implementation of their planners and cannot claim their measured
agent security. AgentDojo offers a foundation for future model-driven
experiments~\cite{agentdojo}, whereas the present suite contains no model.

The current contribution is an artifact and explicit interface contract.
Demonstrating a new cross-modal composition theorem, a reduction in evidence
exchanges or improved useful autonomy against equally expressive controls remains
future research. A common API or a collection of familiar checks is not, by
itself, that scientific contribution.

